
\section{Algoritmos utilizados en SFP}
\label{sec:clasificadores_mt}

La predicción de fallas de software (SFP) es una tarea crítica en la ingeniería de software que busca identificar los módulos del software que son propensos a contener defectos. Para realizar esta predicción, se han utilizado varios algoritmos de aprendizaje automático y técnicas estadísticas. A continuación, se describen algunos de los algoritmos más comúnmente utilizados en la literatura para SFP.

\subsubsection{Regresión Logística}
La regresión logística es un método estadístico utilizado para modelar la probabilidad de que una instancia pertenezca a una de dos clases posibles. Este método ha sido ampliamente utilizado en SFP debido a su simplicidad y efectividad para manejar datos binarios. La regresión logística estima la relación entre las variables independientes (métricas del software) y la variable dependiente (presencia o ausencia de defectos) mediante una función logística \cite{kondo2019impact}.

\subsubsection{Máquinas de Vectores de Soporte (SVM)}
Las máquinas de vectores de soporte (SVM) son un conjunto de métodos de aprendizaje supervisado utilizados para clasificación y regresión. Una SVM busca el hiperplano que separa las dos clases con el máximo margen, es decir, con la mayor distancia posible a las instancias más cercanas de cada clase (los vectores de soporte); la formulación de margen suave permite tolerar instancias mal clasificadas mediante un parámetro de penalización \cite{cortes1995}. En su versión lineal, la SVM es rápida de entrenar sobre conjuntos grandes y sensible a la escala de las características, por lo que requiere estandarizarlas previamente. En el contexto de SFP, la SVM ha demostrado ser eficaz para manejar conjuntos de datos de alta dimensionalidad. En esta tesis se utiliza una SVM lineal con estandarización previa como uno de los clasificadores del análisis de sensibilidad.

\subsubsection{Árboles de Decisión}
Los árboles de decisión son modelos de predicción que utilizan un árbol como estructura de decisión y modelan la relación entre las variables de entrada y la variable de salida. Los árboles de decisión, como el algoritmo J48, han sido utilizados en SFP debido a su capacidad para manejar datos categóricos y continuos, y su facilidad de interpretación \cite{RADJENOVIC20131397}.

\subsubsection{Random Forest}
El algoritmo Random Forest \cite{breiman2001} es un conjunto de árboles de decisión, cada uno entrenado sobre una muestra \textit{bootstrap} de las instancias (\textit{bagging}) y considerando en cada división solo un subconjunto aleatorio de las características. La predicción final se obtiene por votación mayoritaria de los árboles. La combinación de ambas fuentes de aleatoriedad reduce la correlación entre árboles y, con ello, la varianza y el riesgo de sobreajuste del modelo \cite{shah_comparative_2020}. Random Forest es uno de los clasificadores más utilizados en SFP y es el clasificador principal de esta tesis, empleado tanto en Weka \cite{weka} como en \textit{scikit-learn}.

\subsubsection{XGBoost}
XGBoost \cite{chen2016} es una implementación escalable del \textit{boosting} de gradiente sobre árboles de decisión. A diferencia de Random Forest, que entrena árboles independientes y promedia sus votos, el \textit{boosting} construye los árboles de forma secuencial: cada nuevo árbol se ajusta para corregir los errores residuales del conjunto acumulado, minimizando una función de pérdida mediante descenso de gradiente con regularización explícita de la complejidad de los árboles. XGBoost representa, por tanto, un paradigma de ensamble distinto al de Random Forest, y se utiliza en el análisis de sensibilidad de esta tesis para verificar que los hallazgos no son específicos de un único clasificador.

\subsubsection{Naive Bayes}
Naive Bayes es un clasificador basado en la teoría de probabilidad bayesiana. Asume la independencia entre las características, lo cual es una simplificación fuerte pero efectiva en muchos casos. En SFP, Naive Bayes ha sido utilizado debido a su simplicidad y eficiencia computacional, especialmente cuando se dispone de grandes conjuntos de datos con múltiples características.

\subsubsection{K-Nearest Neighbors (KNN)}
El algoritmo K-Nearest Neighbors (KNN) es un método de clasificación que asigna una clase a una instancia basándose en las clases de sus k vecinos más cercanos en el espacio de características. KNN ha sido aplicado en SFP debido a su simplicidad y efectividad para capturar relaciones locales en los datos. Sin embargo, puede ser computacionalmente costoso cuando se trabaja con grandes conjuntos de datos.

\subsubsection{Algoritmos de Sobremuestreo y Submuestreo}
Para abordar el problema del desequilibrio de clases en los conjuntos de datos de SFP, se utilizan técnicas de sobremuestreo y submuestreo. Métodos como SMOTE (Synthetic Minority Over-sampling Technique), ROS (Random Over-Sampling), y RUS (Random Under-Sampling) se emplean para equilibrar las clases defectuosas y no defectuosas, con el objetivo de reducir el sesgo del clasificador hacia la clase mayoritaria \cite{9137899,chawla2002}. Su efecto real, sin embargo, depende del contexto experimental y de la métrica de evaluación \cite{8494821}, un aspecto central de esta tesis.

\subsubsection{Modelos de Conjunto}
Los modelos de conjunto, como el Bagging y el Boosting, combinan múltiples clasificadores para mejorar la precisión y la robustez del modelo. Estas técnicas son útiles en SFP para combinar las predicciones de diferentes algoritmos y obtener una mejor estimación de la propensión a fallos de los módulos de software \cite{9044596}.

En resumen, una variedad de algoritmos de aprendizaje automático y técnicas de preprocesamiento se utilizan en la predicción de fallas de software, cada uno con sus propias ventajas y limitaciones. La selección del algoritmo adecuado depende de las características específicas del conjunto de datos y del problema en cuestión.
